\ifdefined\gkver\else\def\gkver{student}\fi
\documentclass[\gkver]{gaokaozhenti}
\begin{document}

\chapter{1981年全国}

\begin{enumerate}
\item
%% number: 1
%% paperName: 1981年全国高考第 1 题 3 分
%% typeId: 1
%% type: 单选题
%% chapter: 运动和力
%% point: 连接体
%% method: 无
%% score: 3
%% degree: 650
%% duplicateId: 0
%% body:
如图，在光滑的水平桌面上有一物体 A，通过绳子与物体 B 相连。假设绳子的质量以及绳子与定滑轮之间的摩擦力都可忽略不计，绳子不可伸长。如果物体 B 的质量是物体 A 的质量的 3 倍，即 $m_{B}=3m_{A}$，那么物体 A 和 B 的加速度的大小等于\xzanswer{C}
\onepicture[w=4.5cm]{figs/01.png}
\fourchoices[answer=C]
{$3g$}
{$g$}
{$\frac{3}{4}g$}
{$\frac{1}{2}g$}
%% answer: C
%% memo:
\memoanswer{
由题意知，对 A、B 整体进行受力分析，由牛顿第二定律得 $m_{B}g=(m_{A}+m_{B})a$，解得 $a=\frac{3}{4}g$，故 C 正确。
}

\item
%% number: 2
%% paperName: 1981年全国高考第 2 题 3 分
%% typeId: 1
%% type: 单选题
%% chapter: 力物体的平衡
%% point: 摩擦力
%% method: 无
%% score: 3
%% degree: 850
%% duplicateId: 0
%% body:
在光滑的水平桌面上放一物体 A，A 上再放一物体 B，A、B 间有摩擦。施加一水平力 $F$ 于 B，使它相对于桌面向右运动。这时物体 A 相对于桌面\xzanswer{B}
\onepicture[w=4.2cm]{\includesvg{figs/02}}
\fourchoices[answer=B]
{向左动}
{向右动}
{不动}
{运动，但运动方向不能判断}
%% answer: B
%% memo:
\memoanswer{
由于地面光滑，无论 A、B 是否发生相对滑动，A 一定相对桌面向右运动，故 B 正确。

拓展：A 的运动状态有两种可能。

假设 A、B 相对静止一起运动，A、B 间的摩擦力为 $f$，对 A、B 整体，有 $a=\frac{F}{m_{A}+m_{B}}$；

对 A，有 $f=m_{A}a=\frac{m_{A}F}{m_{A}+m_{B}}$。

又 $f_{m}=\mu m_{B}g$。

A、B 相对静止需满足 $f\le f_{m}=\mu m_{B}g$，联立解得 $\mu\ge\frac{m_{A}F}{m_{B}(m_{B}+m_{A})}$。

若 $\mu\ge\frac{m_{A}F}{m_{B}(m_{B}+m_{A})}$，A、B 一起向右运动；

若 $\mu<\frac{m_{A}F}{m_{B}(m_{B}+m_{A})}$，A、B 发生相对滑动，均向右运动，但 B 运动得更快。
}

\item
%% number: 3
%% paperName: 1981年全国高考第 3 题 3 分
%% typeId: 1
%% type: 单选题
%% chapter: 电路
%% point: 电容
%% method: 无
%% score: 3
%% degree: 850
%% duplicateId: 0
%% body:
平行板电容器，其两板始终保持和一直流电源的正、负极相连接，当两板间插入电介质时，电容器的带电量和两板间的电势差的变化是\xzanswer{C}
\fourchoices[answer=C]
{带电量不变，电势差增大}
{带电量不变，电势差减小}
{带电量增大，电势差不变}
{带电量减小，电势差不变}
%% answer: C
%% memo:
\memoanswer{
由平行板电容器的决定式 $C=\frac{\varepsilon S}{4\pi kd}$ 知，平行板电容器中插入电介质后，介电常数 $\varepsilon$ 会增大而其他量不变，故电容增大；由题意知两板始终与电源连接，故电压不变，由 $Q=UC$ 知 $U$ 不变、$C$ 增大，所以 $Q$ 增大，故 C 正确。
}

\item
%% number: 4
%% paperName: 1981年全国高考第 4 题 3 分
%% typeId: 1
%% type: 单选题
%% chapter: 电场
%% point: 静电感应
%% method: 无
%% score: 3
%% degree: 550
%% duplicateId: 0
%% body:
把一个架在绝缘支座上的导体放在负电荷形成的电场中，导体处于静电平衡时，导体表面上感应电荷的分布如图所示，这时导体\xzanswer{D}
\onepicture[w=4.6cm]{figs/04.png}
\fourchoices[answer=D]
{A 端的电势比 B 端的电势高}
{A 端的电势比 B 端的电势低}
{A 端的电势可能比 B 端的电势高，也可能比 B 端的电势低}
{A 端的电势与 B 端的电势相等}
%% answer: D
%% memo:
\memoanswer{
导体在点电荷附近，出现静电感应现象，导致导体表面的电子重新分布，因此在导体内部出现感应电荷的电场，正好与点电荷的电场叠加；静电平衡时内部电场强度处处为零，导体中的电子不再移动，故导体是等势体，表面是等势面，即 A 端的电势与 B 端的电势相等，故 D 正确。
}

\item
%% number: 5
%% paperName: 1981年全国高考第 5 题 3 分
%% typeId: 1
%% type: 单选题
%% chapter: 电路
%% point: 电阻定律
%% method: 无
%% score: 3
%% degree: 750
%% duplicateId: 0
%% body:
一段粗细均匀的镍铬丝，横截面的直径是 $d$，电阻是 $R$。把它拉制成直径是 $\frac{1}{10}d$ 的均匀细丝后，它的电阻变成\xzanswer{B}
\fourchoices[answer=B]
{$\frac{1}{10000}R$}
{$10000R$}
{$\frac{1}{100}R$}
{$100R$}
%% answer: B
%% memo:
\memoanswer{
当镍铬丝的直径为 $d'=\frac{d}{10}$ 时，则截面积为 $S'=\frac{S}{100}$；由于镍铬丝体积不变，其长度必相应地变为 $L'=100L$，由电阻定律得 $R'=\rho\frac{L'}{S'}=\rho\frac{100L}{\frac{1}{100}S}=10000\rho\frac{L}{S}=10000R$，故 B 正确。
}

\item
%% number: 6
%% paperName: 1981年全国高考第 6 题 3 分
%% typeId: 1
%% type: 单选题
%% chapter: 曲线运动
%% point: 万有引力定律
%% method: 无
%% score: 3
%% degree: 550
%% duplicateId: 0
%% body:
假设火星和地球都是球体，火星的质量 $M_{\text{火}}$ 和地球的质量 $M_{\text{地}}$ 之比 $M_{\text{火}}/M_{\text{地}}=p$，火星的半径 $R_{\text{火}}$ 和地球的半径 $R_{\text{地}}$ 之比 $R_{\text{火}}/R_{\text{地}}=q$，那么火星表面处的重力加速度 $g_{\text{火}}$ 和地球表面处的重力加速度 $g_{\text{地}}$ 之比 $g_{\text{火}}/g_{\text{地}}$ 等于\xzanswer{A}
\fourchoices[answer=A]
{$p/q^{2}$}
{$pq^{2}$}
{$p/q$}
{$pq$}
%% answer: A
%% memo:
\memoanswer{
星球表面的物体受到的重力等于万有引力，有 $G\frac{Mm}{R^{2}}=mg$，解得 $g=\frac{GM}{R^{2}}$，所以
\[ \frac{g_{\text{火}}}{g_{\text{地}}}=\frac{M_{\text{火}}}{M_{\text{地}}}\cdot\frac{R_{\text{地}}^{2}}{R_{\text{火}}^{2}}=\frac{p}{q^{2}}, \]
故 A 正确。
}

\item
%% number: 7
%% paperName: 1981年全国高考第 7 题 4 分
%% typeId: 3
%% type: 填空题
%% chapter: 气体性质
%% point: 气体的状态参量
%% method: 无
%% score: 4
%% degree: 950
%% duplicateId: 0
%% body:
1 标准大气压 = \tkanswer{760} 毫米水银柱 = \tkanswer{$1.01\times10^{5}$} 帕斯卡（即牛顿/米$^{2}$）（水银的密度 $\rho=13.6$ 克/厘米$^{3}$，重力加速度 $g=9.81$ 米/秒$^{2}$。本题答案要求取三位有效数字）。
%% answer: 760；$1.01\times10^{5}$
\jdanswer{
760；$1.01\times10^{5}$
}
%% memo:
\memoanswer{
1 标准大气压 $=760\UmmHg=10.339\,\mathrm{mH_{2}O}=1.01325\times10^{5}\UPa$。

评分说明：共 4 分。第一个答案正确的给 2 分，错误的给 0 分；第二个答案，正确的给 2 分；数量级和前两位有效数字正确而第三位有效数字有出入的，给 1 分；数量级不对的，给 0 分。
}

\item
%% number: 8
%% paperName: 1981年全国高考第 8 题 4 分
%% typeId: 3
%% type: 填空题
%% chapter: 光学
%% point: 透镜
%% method: 无
%% score: 4
%% degree: 850
%% duplicateId: 0
%% body:
一照相机，用焦距 $f=0.20$ 米的凸透镜做镜头，用来为一个站立在镜头前 4.2 米处的儿童照相时，底片应该离镜头 \tkanswer{0.21} 米，底片上像的高度和儿童的身高之比是 1: \tkanswer{20}。
%% answer: 0.21；20
\jdanswer{
0.21；20
}
%% memo:
\memoanswer{
由成像公式得 $\frac{1}{f}=\frac{1}{v}+\frac{1}{u}$，由题意知 $f=0.20\Um$，$u=4.2\Um$，解得 $v=0.21\Um$；由放大率公式得 $M=\frac{v}{u}=\frac{0.21}{4.2}=\frac{1}{20}$，故底片上的像与儿童的身高之比为 1:20。

评分标准：共 4 分。两个答案都要求准确到第二位有效数字。每个答案正确的给 2 分，错误的给 0 分。
}

\item
%% number: 9
%% paperName: 1981年全国高考第 9 题 4 分
%% typeId: 3
%% type: 填空题
%% chapter: 电路
%% point: 电功电功率
%% method: 无
%% score: 4
%% degree: 650
%% duplicateId: 0
%% body:
把 5 欧姆的电阻 $R_{1}$ 和 10 欧姆的电阻 $R_{2}$ 串联起来，然后在这段串联电路的两端加 15 伏特的电压，这时 $R_{1}$ 消耗的电功率是 \tkanswer{5} 瓦特，$R_{2}$ 消耗的电功率是 \tkanswer{10} 瓦特。把 $R_{1}$ 和 $R_{2}$ 改为并联，如果要使 $R_{1}$ 仍消耗与原来同样大小的电功率，则应在它们两端加 \tkanswer{5} 伏特的电压，这时 $R_{2}$ 消耗的电功率是 \tkanswer{2.5} 瓦特。
%% answer: $5\UW$；$10\UW$；$5\UV$；$2.5\UW$
\jdanswer{
$5\UW$；$10\UW$；$5\UV$；$2.5\UW$
}
%% memo:
\memoanswer{
（1）两电阻串联，两端加 $15\UV$ 电压时，电路中电流为
\[ I=\frac{U}{R_{1}+R_{2}}=\frac{15\UV}{5\UO+10\UO}=1\UA, \]
电阻消耗的电功率为
\[ P_{1}=I^{2}R_{1}=(1\UA)^{2}\times5\UO=5\UW, \]
\[ P_{2}=I^{2}R_{2}=(1\UA)^{2}\times10\UO=10\UW. \]
又 $P=\frac{U^{2}}{R}$。

（2）两电阻并联时，两端电压为 $U'$，电阻消耗的电功率为
\[ P_{1}'=\frac{U'^{2}}{R_{1}}=\frac{U'^{2}}{5}, \]
又 $P_{1}=P_{1}'$，解得 $U'=5\UV$，此时电阻 $R_{2}$ 消耗的电功率为
\[ P_{2}'=\frac{U'^{2}}{R_{2}}=\frac{(5\UV)^{2}}{10\UO}=2.5\UW. \]

评分说明：共 4 分，每个答案正确的，给 1 分，错误的，给 0 分。
}

\item
%% number: 10
%% paperName: 1981年全国高考第 10 题 4 分
%% typeId: 3
%% type: 填空题
%% chapter: 机械波
%% point: 波的图象
%% method: 无
%% score: 4
%% degree: 850
%% duplicateId: 0
%% body:
右图是一列沿 $x$ 轴正方向传播的机械横波在某一时刻的图像。从图上可看出，这列波的振幅是 \tkanswer{0.03} 米，波长是 \tkanswer{2} 米，P 处的质点在此时刻的运动方向 \tkanswer{沿 $y$ 轴正方向（或向上）}。
\onepicture[align=right, w=6cm]{figs/10.png}
%% answer: 0.03；2；沿 $y$ 轴正方向（或向上）
\jdanswer{
0.03；2；沿 $y$ 轴正方向（或向上）
}
%% memo:
\memoanswer{
由图可知，振幅 $A=0.03\Um$，波长 $\lambda=2\Um$，利用同侧法可判断，此时的 P 竖直向上振动，即沿 $y$ 轴正方向（或向上）运动。

评分说明：共 4 分。前两个答案，正确的各给 1 分，错误的给 0 分；第三个答案正确的给 2 分；未用文字回答而画一向上箭头表示方向的，同样给 2 分；错误的给 0 分。
}

\item
%% number: 11
%% paperName: 1981年全国高考第 11 题 4 分
%% typeId: 3
%% type: 填空题
%% chapter: 运动和力
%% point: 动量和冲量
%% method: 无
%% score: 4
%% degree: 850
%% duplicateId: 0
%% body:
质量是 $m$ 的质点，以匀速率 $v$ 作圆周运动，圆心在坐标系的原点 $O$。在质点从位置 1 运动到位置 2（如右图所示）的过程中，作用在质点上的合力的功等于 \tkanswer{0}；合力冲量的大小是 \tkanswer{$\sqrt{2}mv$}，方向与 $x$ 轴正方向成 \tkanswer{$225^{\circ}$}（逆时针计算角度）。
\onepicture[align=right, w=3.4cm]{\includesvg{figs/11}}
%% answer: 0；$\sqrt{2}mv$；$225^{\circ}$
\jdanswer{
0；$\sqrt{2}mv$；$225^{\circ}$
}
%% memo:
\memoanswer{
质点由 1 到 2 的过程中，由动能定理得 $W_{F_{\text{合}}}=\Delta E_{\text{k}}$；因物体做匀速圆周运动，所以 $\Delta E_{\text{k}}=0$，故 $W_{F_{\text{合}}}=0$。由动量定理可知 $I_{F_{\text{合}}}=\Delta p$，质点在位置 1 时，动量为 $mv$，方向竖直向上；质点在位置 2 时，动量为 $mv$，方向水平向左，由矢量三角形定则可知 $\Delta p=\sqrt{2}mv$，方向与 $x$ 轴正方向成 $225^{\circ}$（或 $\frac{5}{4}\pi$）（逆时针计算角度）。

评分说明：共 4 分。第一个答案 1 分；第二个答案，正确的给 2 分，错误的给 0 分；第三个答案，正确的给 1 分。
}

\item
%% number: 12
%% paperName: 1981年全国高考第 12 题 4 分
%% typeId: 3
%% type: 填空题
%% chapter: 原子和原子核
%% point: 人工核反应
%% method: 无
%% score: 4
%% degree: 850
%% duplicateId: 0
%% body:
铝核 $\ce{^{27}_{13}Al}$ 被 $\alpha$ 粒子击中后产生的反应生成物是磷核 $\ce{^{30}_{15}P}$，同时放出一个 \tkanswer{中子（或 $\ce{^{1}_{0}n}$）}，这个核反应方程是 \tkanswer[6cm]{$\ce{^{27}_{13}Al + ^{4}_{2}He -> ^{30}_{15}P + ^{1}_{0}n}$}。
%% answer: 中子（或 $\ce{^{1}_{0}n}$）；$\ce{^{27}_{13}Al + ^{4}_{2}He -> ^{30}_{15}P + ^{1}_{0}n}$
\jdanswer{
中子（或 $\ce{^{1}_{0}n}$）；$\ce{^{27}_{13}Al + ^{4}_{2}He -> ^{30}_{15}P + ^{1}_{0}n}$
}
%% memo:
\memoanswer{
核反应方程满足核电荷数和质量数守恒。由核电荷数守恒得 $Z=13+2-15=0$，该粒子的核电荷数为 0；由质量数守恒得 $A=27+4-30=1$，该粒子的质量数为 1，故该粒子为中子（或 $\ce{^{1}_{0}n}$），核反应方程为 $\ce{^{27}_{13}Al + ^{4}_{2}He -> ^{30}_{15}P + ^{1}_{0}n}$。

评分说明：共 4 分。每个答案，正确的给 2 分，错误的给 0 分。
}

\item
%% number: 13
%% paperName: 1981年全国高考第 13 题 3 分
%% typeId: 4
%% type: 实验题
%% chapter: 直线运动
%% point: 游标卡尺和螺旋测微仪
%% method: 无
%% score: 3
%% degree: 950
%% duplicateId: 0
%% body:
用游标卡尺（图~\ref{1981全国13a}）测一根金属管的内径和外径时，卡尺上的游标位置分别如图~\ref{1981全国13b}和图~\ref{1981全国13c}所示。这根金属管的内径读数是 \tkanswer{2.37} 厘米，外径读数是 \tkanswer{3.03} 厘米，管壁厚是 \tkanswer{0.33} 厘米。
\threepicture[numstyle=arabic, labelA=1981全国13a, labelB=1981全国13b, labelC=1981全国13c, wA=7cm, wB=4.5cm, wC=4.5cm, align=right]
{figs/13a.png}
{figs/13b.png}
{figs/13c.png}
%% answer: 2.37；3.03；0.33
\jdanswer{
2.37；3.03；0.33
}
%% memo:
\memoanswer{
图 2：主尺读数 $23\Umm$，游标上的第 7 个小分格与主尺上的刻度恰好对正，所以内径读数为 $23\Umm+7\times0.01\Umm=23.7\Umm=2.37\Ucm$；同理图 3：外径的读数为 $30\Umm+3\times0.01\Umm=30.3\Umm=3.03\Ucm$；管壁厚应等于外径与内径差的一半，即 $\frac{3.03\Ucm-2.37\Ucm}{2}=0.33\Ucm$。

评分说明：共 3 分。每个答案，正确的给 1 分，错误的给 0 分。
}

\item
%% number: 14
%% paperName: 1981年全国高考第 14 题 4 分
%% typeId: 4
%% type: 实验题
%% chapter: 力物体的平衡
%% point: 重力
%% method: 无
%% score: 4
%% degree: 850
%% duplicateId: 0
%% body:
用下图所示的天平称质量前，先要进行哪些调节？说明调节哪些部件和怎样才算调节好了。
\onepicture[align=right, w=5.5cm]{figs/14.png}
%% answer: 要进行两步调节：1．使天平的底板 B 水平。调节螺旋 S，直到重垂线 Q 的小锤尖端跟小锥体Z的尖端对正，这就表示底板水平了。2．使天平平衡。调节螺旋 S$'$，使指针 D 指在标尺 K 的中央，这就表示天平平衡了。
\jdanswer{
要进行两步调节：（1）使天平的底板 B 水平。调节螺旋 S，直到重垂线 Q 的小锤尖端跟小锥体 Z 的尖端对正，这就表示底板水平了。（2）使天平平衡。调节螺旋 $S'$，使指针 D 指在标尺 K 的中央，这就表示天平平衡了。
}
%% memo:
\memoanswer{
要进行两步调节：（1）使天平的底板 B 水平。调节螺旋 S，直到重垂线 Q 的小锤尖端跟小锥体 Z 的尖端对正，这就表示底板水平了。（2）使天平平衡。调节螺旋 $S'$，使指针 D 指在标尺 K 的中央，这就表示天平平衡了。

评分说明：共 4 分。第一步调节，占 2 分；只答出要调节底板水平的，给 1 分；既答对调节 S 又答出调节好了的标志的，再给 1 分，二者缺一的，不再给分。第二步调节，也占 2 分；只答出要调节天平平衡的，给 1 分；既答对调节 $S'$ 又答出调节好了的标志的，再给 1 分，二者缺一的，不再给分。
}

\item
%% number: 15
%% paperName: 1981年全国高考第 15 题 3 分
%% typeId: 4
%% type: 实验题
%% chapter: 电路
%% point: 多用表的使用
%% method: 无
%% score: 3
%% degree: 750
%% duplicateId: 0
%% body:
用万用电表电阻挡判断一只 PNP 型晶体三极管的基极时，电表指针的偏转情况如下图所示。哪只管脚是基极？
\onepicture[align=right, w=10cm]{figs/15.png}
%% answer: 管脚 3 是基极
\jdanswer{
管脚 3 是基极
}
%% memo:
\memoanswer{
管脚 3 是基极。

评分说明：本小题 3 分。用文字答出管脚 3 是基极，或在图中管脚 3 旁注明基极（或注明 b）的，都给 3 分；不要求说明理由。
}

\item
%% number: 16
%% paperName: 1981年全国高考第 16 题 10 分
%% typeId: 6
%% type: 计算题
%% chapter: 气体性质
%% point: 气体多过程
%% method: 无
%% score: 10
%% degree: 650
%% duplicateId: 0
%% body:
使一定质量的理想气体的状态按图~\ref{1981全国16a}中箭头所示的顺序变化，图线 BC 是一段以纵轴和横轴为渐近线的双曲线。
\begin{enumerate}
	\item
	已知气体在状态 A 的温度 $T_{A}=300\UK$，求气体在状态 B、C 和 D 的温度各是多少。
	\item
	将上述状态变化过程在图~\ref{1981全国16b}中画成用体积 $V$ 和温度 $T$ 表示的图线（图中要标明 A、B、C、D 四点，并且要画箭头表示变化的方向）。说明每段图线各表示什么过程。
\end{enumerate}
\twopicture[numstyle=arabic, labelA=1981全国16a, labelB=1981全国16b, wA=5.2cm, wB=5.4cm, align=right]
{figs/16a.png}
{figs/16b.png}
%% answer: （1）$T_{B}=T_{C}=600\UK$，$T_{D}=300\UK$；（2）见解析图
\jdanswer{
\begin{enumerate}
	\item
	$T_{B}=T_{C}=600\UK$，$T_{D}=300\UK$

	\item
	见解析图

\end{enumerate}
}
%% memo:
\memoanswer{
（1）设标准大气压的压强为 $p_{0}$，由 $p-V$ 图象可知，气体在 A、B、C、D 各状态下的压强和体积分别为 $p_{A}=4p_{0}$、$V_{A}=10\UL$；$p_{B}=4p_{0}$、$V_{B}=20\UL$；$p_{C}=2p_{0}$、$V_{C}=40\UL$；$p_{D}=2p_{0}$、$V_{D}=20\UL$。

已知 $T_{A}=300\UK$，设气体在 C、D 各状态下的温度分别为 $T_{C}$、$T_{D}$，又 BC 为渐近线（双曲线），故 $T_{B}=T_{C}$，由理想气体状态方程得
\[ \frac{p_{A}V_{A}}{T_{A}}=\frac{p_{B}V_{B}}{T_{B}}=\frac{p_{C}V_{C}}{T_{C}}=\frac{p_{D}V_{D}}{T_{D}}, \]
即 $\frac{4p_{0}\times10}{300}=\frac{4p_{0}\times20}{T_{B}}=\frac{2p_{0}\times40}{T_{C}}=\frac{2p_{0}\times20}{T_{D}}$，
解得 $T_{C}=T_{B}=600\UK$，$T_{D}=300\UK$。

（2）在 $V-T$ 图上状态变化过程的图线如下：

\onepicture[w=5.2cm]{figs/16_an.png}

AB 是等压过程，BC 是等温过程，CD 是等压过程。

评分说明：全题 10 分。（1）3 分，（2）7 分。（1）中，每个答案，正确的给 1 分；数值错误的，给 0 分；数值正确而缺单位或单位错误的，无论在几个答案中出现，只扣 1 分。直接由 $p-V$ 图上断定 $V_{B}=20\UL$，并正确求出 $T_{B}$ 的，同样给分。（2）中，A、B、C、D 四个点，每个点正确标明的，各给 1 分；AB、BC、CD 三条图线画正确并说明是什么过程的，各给 1 分；图线画对而未说明是什么过程的，各给 0 分；在图上未画箭头或三个箭头未画全的，扣 1 分；图上两个坐标轴每有一个未注明坐标分度值的，扣 1 分。
}

\item
%% number: 17
%% paperName: 1981年全国高考第 17 题 10 分
%% typeId: 6
%% type: 计算题
%% chapter: 光学
%% point: 光电效应
%% method: 无
%% score: 10
%% degree: 650
%% duplicateId: 0
%% body:
一光电管的阴极用极限波长 $\lambda_{0}=5000$ 埃的钠制成。用波长 $\lambda=3000$ 埃的紫外线照射阴极，光电管阳极 A 和阴极 K 之间的电势差 $U=2.1\UV$，光电流的饱和值 $I=0.56\UmuA$。
\begin{enumerate}
	\item
	求每秒内由 K 极发射的电子数；
	\item
	求电子到达 A 极时的最大动能；
	\item
	如果电势差 $U$ 不变，而照射光的强度增到原值的三倍，此时电子到达 A 极时的最大动能是多大？
\end{enumerate}

（普朗克恒量 $h=6.63\times10^{-34}\UJcdots$，电子电量 $e=1.60\times10^{-19}\UC$，真空中的光速 $c=3.00\times10^{8}\Ums$）
\onepicture[align=right, w=4cm]{figs/17.png}
%% answer: （1）$3.5\times10^{12}$；（2）$6.01\times10^{-19}\UJ$；（3）以上结果不变
\jdanswer{
\begin{enumerate}
	\item
	$3.5\times10^{12}$

	\item
	$6.01\times10^{-19}\UJ$

	\item
	以上结果不变

\end{enumerate}
}
%% memo:
\memoanswer{
（1）每秒内发射的电子数为
\[ n=\frac{It}{e}=\frac{0.56\times10^{-6}\UA\times1\Us}{1.60\times10^{-19}\UC}=3.5\times10^{12}. \]（a）

（2）每个紫外线光子的能量为
\[ E_{\text{紫}}=h\nu=h\frac{c}{\lambda}=6.63\times10^{-19}\UJ, \]（b）
钠的逸出功为
\[ W_{0}=\frac{hc}{\lambda_{0}}=3.98\times10^{-19}\UJ, \]（c）
每个电子在电场中被加速而获得的能量为
\[ eU=1.60\times10^{-19}\UC\times2.1\UV=3.36\times10^{-19}\UJ, \]（d）
由能量守恒定律得，电子的最大动能为
\[ \frac{1}{2}mv^{2}=\frac{hc}{\lambda}-\frac{hc}{\lambda_{0}}+eU=6.01\times10^{-19}\UJ=3.76\UeV. \]（e）

（3）最大动能与入射光的频率有关，与光强无关，当光强增到 3 倍时，最大初动能不变，为 $6.01\times10^{-19}\UJ$（或 $=3.76\UeV$）。

评分说明：全题 10 分。（1）3 分，（2）6 分，（3）1 分。（1）中，正确算出结果（a）的，给 3 分；单纯数字计算错误（包括数字正确而数量级错误）扣 1 分。（2）中，正确算出光子能量（b）的，给 1 分；正确算出逸出功（c）的，再给 2 分；正确算出电子在电场中获得的能量（d）的，再给 1 分；根据能量守恒关系正确算出电子最大动能的，再给 2 分。直接列出式（e）并算出正确结果的，同样给 6 分。如果考生根据爱因斯坦方程，正确算出电子离开 K 时的最大动能，即只漏去 $eU$ 一项，把它当作本题答案的，给 4 分；单纯数字计算错误扣 1 分；数值正确而缺单位或单位错误的，扣 1 分。（3）中，答案正确的，给 1 分。
}

\item
%% number: 18
%% paperName: 1981年全国高考第 18 题 14 分
%% typeId: 6
%% type: 计算题
%% chapter: 电磁感应
%% point: 切割磁感线电动势
%% method: 无
%% score: 14
%% degree: 550
%% duplicateId: 0
%% body:
用均匀导线弯成正方形闭合线框 abcd，线框每边长 $8.0\Ucm$，每边的电阻值为 $0.010\UO$。把线框放在磁感应强度为 $B=0.050\UT$ 的匀强磁场中，并使它绕轴 $OO'$ 以 $\omega=100\Urads$ 的匀角速度旋转，旋转方向如图所示。已知轴 $OO'$ 在线框平面内，并且垂直于 $B$，$Od=3Oa$，$O'c=3O'b$。当线框平面转至和 $B$ 平行的瞬时（如图所示），
\begin{enumerate}
	\item
	每个边产生的感生电动势的大小各是多少？
	\item
	线框内感生电流的大小是多少？在图中用箭头标出感生电流的方向；
	\item
	e、f 分别为 ab 和 cd 的中点，e、f 两点间的电势差 $U_{ef}$（即 $U_{e}-U_{f}$）是多大？
\end{enumerate}
\onepicture[align=right, w=4.4cm]{\includesvg{figs/18}}
%% answer: （1）cd 段感生电动势的大小为 $\varepsilon_{1}=0.024\UV$，ab 段感生电动势的大小为 $\varepsilon_{2}=0.008\UV$，da 段和 bc 段的电动势都是零；（2）$I=0.8\UA$，电流方向沿 dcbad；（3）$U_{ef}=0\UV$
\jdanswer{
\begin{enumerate}
	\item
	cd 段感生电动势的大小为 $\varepsilon_{1}=0.024\UV$，ab 段感生电动势的大小为 $\varepsilon_{2}=0.008\UV$，da 段和 bc 段的电动势都是零

	\item
	$I=0.8\UA$，电流方向沿 dcbad。

	\item
	$U_{ef}=0\UV$

\end{enumerate}
}
%% memo:
\memoanswer{
（1）令 $l$ 表示每边边长，$R$ 表示其电阻值，则 $l=8.0\Ucm$，$R=0.010\UO$。设 cd 段感应电动势的大小为 $E_{1}$，ab 段感应电动势的大小为 $E_{2}$，则
\[ E_{1}=Blv_{1}\sin90^{\circ}=Bl\omega\cdot\frac{3}{4}l=\frac{3}{4}Bl^{2}\omega=0.024\UV, \]①
\[ E_{2}=Blv_{2}\sin90^{\circ}=Bl\omega\cdot\frac{1}{4}l=\frac{1}{4}Bl^{2}\omega=0.008\UV, \]②
da 段和 bc 段不切割磁感线，所以它们的电动势都是零。

（2）线框中的总电动势为
\[ E=E_{1}+E_{2}=0.032\UV, \]①
线框中的感应电流为
\[ I=\frac{E}{4R}=\frac{0.032}{0.040}\UA=0.80\UA, \]②
由楞次定律或右手定则，可判断电流方向沿 dcbad。

（3）解法一：$\varphi_{e}-\varphi_{f}$ 是 ebcf 一段电路两端的电势差（逆时针观看），它应等于 eb 段的路端电压 $U_{eb}$、bc 段两端的电势差 $U_{bc}$ 与 cf 段的路端电压 $U_{cf}$ 的代数和，即
\[ U_{ef}=U_{eb}+U_{bc}+U_{cf}, \]①
\[ U_{eb}=\frac{1}{2}E_{2}-I\frac{R}{2}=0.004\UV-0.8\UA\times0.005\UO=0\UV, \]②
\[ U_{bc}=-IR=-0.8\UA\times0.010\UO=-0.008\UV, \]③
\[ U_{cf}=\frac{1}{2}E_{1}-I\frac{R}{2}=0.012\UV-0.8\UA\times0.005\UO=0.008\UV, \]④
故 $U_{ef}=U_{eb}+U_{bc}+U_{cf}=0\UV$。

解法二：$\varphi_{e}-\varphi_{f}$ 是 eadf 一段电路两端的电势差，根据题意，由电势叠加得
\[ \varphi_{e}+\frac{1}{2}E_{2}-I\left(\frac{R}{2}+R+\frac{R}{2}\right)+\frac{1}{2}E_{1}=\varphi_{f}, \]
又 $U_{ef}=\varphi_{e}-\varphi_{f}$，联立解得 $U_{ef}=0\UV$。

解法三：闭合电路 dcbad 中的总电动势等于总电势降落。在 ebcf 一段电路中，电动势等于总电动势的一半；电阻等于总电阻的一半，因而电势降落是总电势降落的一半，于是，在这段电路中，电势升正好等于电势降，由此可见，两端的电势差等于零。

评分说明：全题 14 分。（1）5 分，（2）4 分，（3）5 分。（1）中，正确求出 $E_{1}$ 和 $E_{2}$ 的，各给 2 分；答出 da 和 bc 段电动势都是零的，合给 1 分；单纯数字计算错误扣 1 分；$E_{1}$ 和 $E_{2}$ 的答案数值正确而缺单位或单位错误的，无论出现一次或二次，只扣 1 分。因把 $v=\omega r$ 的关系搞错而引起答案错误的，只扣 1 分。（2）中，正确算出总电动势的，给 1 分；进一步正确算出电流的，再给 2 分；直接求出结果（d）的，同样给 3 分；电流方向正确的再给 1 分；数值错误和单位错误的扣分同（1）中规定。（3）中，（e）、（f）、（g）和（h）各占 1 分；利用以上四式进一步正确算出结论（i）的，再给 1 分；不分步计算，直接正确列出（i）式并算出结果的，同样给 5 分；单纯运算错误的，扣 1 分。只有 $U_{ef}$ 等于零的结论，而无任何推算过程、无任何论述或论述错误的，均不给分；结论正确但论述不够清楚的，酌情给分。
}

\item
%% number: 19
%% paperName: 1981年全国高考第 19 题 14 分
%% typeId: 6
%% type: 计算题
%% chapter: 机械振动
%% point: 弹簧振子
%% method: 无
%% score: 14
%% degree: 350
%% duplicateId: 0
%% body:
在光滑水平面的两端对立着两堵竖直的墙 A 和 B，把一根倔强系数是 $k$ 的弹簧的左端固定在墙 A 上，在弹簧右端系一个质量是 $m$ 的物体 1。用外力压缩弹簧（在弹性限度内）使物体 1 从平衡位置 O 向左移动距离 $s_{0}$，紧靠着 1 放一个质量也是 $m$ 的物体 2，使弹簧、1 和 2 都处于静止状态，然后撤去外力，由于弹簧的作用，物体开始向右滑动。
\begin{enumerate}
	\item
	在什么位置物体 2 与物体 1 分离？分离时物体 2 的速率是多大？
	\item
	物体 2 离开物体 1 后继续向右滑动，与墙 B 发生完全弹性碰撞。B 与 O 之间的距离 $x$ 应满足什么条件，才能使 2 在返回时恰好在 O 点与 1 相遇？设弹簧的质量以及 1 和 2 的宽度都可忽略不计。
\end{enumerate}
\onepicture[align=right, w=11cm]{figs/19.png}
%% answer: （1）在 O 点分离，$v=\sqrt{\frac{k}{2m}}s_{0}$；（2）$x=n\frac{\sqrt{2}\pi}{4}s_{0}$，$n=1,2,3\ldots$
\jdanswer{
\begin{enumerate}
	\item
	到达平衡位置 O 前，1 和 2 一起做加速运动；到 O 点后，1 开始减速，2 开始做匀速运动，因而 2 和 1 将在 O 点分离，分离时 $v=\sqrt{\frac{k}{2m}}s_{0}$

	\item
	$x=n\frac{\sqrt{2}\pi}{4}s_{0}$，$n=1,2,3\ldots$

\end{enumerate}
}
%% memo:
\memoanswer{
（1）到达平衡位置 O 前，1 和 2 一起做加速运动，到 O 点后，1 开始减速，2 开始做匀速运动，因而 2 和 1 将在 O 点分离。

到达 O 点前，把 1、2 和弹簧看做一个系统，只有系统内的弹簧的弹性力做功，所以系统的机械能守恒。令 $v$ 表示 1 和 2 到达 O 点时的速率，则有
\[ \frac{1}{2}(m+m)v^{2}=\frac{1}{2}ks_{0}^{2}, \]
解得 $v=\frac{\sqrt{2}}{2}\sqrt{\frac{k}{m}}s_{0}$，故分离时物体 2 的速率为 $v=\sqrt{\frac{k}{2m}}s_{0}$。

（2）分离后，在下一次相遇前，1 以 O 点为平衡位置做简谐振动，振动的周期为
\[ T=2\pi\sqrt{\frac{m}{k}}, \]①
从 1 和 2 分离时开始计时，即令该时刻 $t=0$，则 1 通过 O 点的时刻为
\[ t_{1}=\frac{1}{2}T,\ T,\ \frac{3}{2}T,\ 2T,\ \ldots=n\cdot\frac{1}{2}T,\quad n=1,2,3,\ldots, \]②
过 O 点后，2 以匀速率 $v$ 向右做直线运动，与 B 相碰时，由于碰撞是完全弹性的，碰撞后 2 的速率不变，运动反向。

令 $x$ 表示 B 与 O 点间的距离，则 2 返回 O 点的时刻为
\[ t_{2}=\frac{2x}{v}, \]③
如 2 恰好在 O 点与 1 相遇，则
\[ t_{2}=t_{1}, \]④
将②③两式代入④，即得 $x$ 应满足的条件为
\[ x=n\frac{\sqrt{2}\pi}{4}s_{0},\quad n=1,2,3,\ldots \]⑤

评分说明：全题 14 分。（1）4 分，（2）10 分。（1）中，答出在 O 点分离的，给 2 分；列出机械能守恒方程并求出 2 的速率的，再给 2 分。（2）中，知道和 2 分离后 1 作简谐振动，并写出振动周期公式（a）的，给 2 分；正确列出 1 经过 O 点的时刻 $t_{1}$，即式（b）的，再给 4 分；对于 $t_{1}$，只回答了 $n=1$ 或 $n=2$ 一次的，扣 3 分；只回答了 $n=1$ 和 $n=2$ 两次的，扣 2 分；只回答了 $n$ 为奇数或为偶数一种情形的，扣 2 分；由于这一步考虑不全面，导致本题最后答案不全的，后面不重复扣分。答出 2 与墙 B 碰撞后，速率不变，运动反向的（不要求证明），给 1 分，又正确求出 2 返回 O 点的时刻（c）的，再给 1 分。正确列出 1 和 2 在 O 点相遇的条件，即（d）的，给 1 分；进一步求出距离 $x$，即（e）的，再给 1 分。将（d）、（e）两步并作一步直接求出结果的，同样给 2 分。本题中的单纯演算错误，可视其对物理过程或最后结果的影响程度，酌情扣分。
}

\end{enumerate}

\end{document}
